\providecommand{\cH}{\mathcal{H}}
\providecommand{\cB}{\mathcal{B}}
\providecommand{\Barch}{\overline{B}^{\mathrm{ch}}}

% Master synthesis package for the Yang--Mills / bar-cobar / W-boundary / instanton / IR program.
% This file is designed as a single integration-ready module collecting and harmonizing all
% theorem packages produced in the current thread.
% Suggested insertion point: after Chapter 31 or as a dedicated appendix / research note.

\section{Toward the platonic ideal: twisted and untwisted Yang--Mills through bar--cobar, higher operations, instantons, and the infrared}
% label removed: sec:ym-platonic-ideal-synthesis

\begin{summary}[Deep essential]
The whole programme can be organized around a single chain of ideas:
\[
\text{boundary chiral algebra}
\longrightarrow
\text{Koszul dual / derived center}
\longrightarrow
\text{mixed and higher-body couplings}
\longrightarrow
\text{instanton-completed visible center}
\longrightarrow
\text{visible defect category}
\longrightarrow
\text{screening Laplacian}
\longrightarrow
\text{conditional mass gap}.
\]
The first four arrows are internal to the mathematics below.
The passage from the visible defect category to untwisted Yang--Mills dynamics remains conditional:
it reduces to constructing a natural untwisting bridge datum and verifying a screening domination estimate.
\end{summary}

\begin{theorem}[Grand synthesis principle; \ClaimStatusConditional]
% label removed: thm:grand-synthesis-principle
Fix a smooth projective curve $X$ and a screened instanton-completed boundary simplex built from
Koszul-admissible twisted Yang--Mills boundary data.
Then:
\begin{enumerate}[label=(\roman*)]
\item every first-order boundary coupling is dual to a class in the instanton-completed derived center;
\item every genuine higher-body visible coupling is dual to a higher \v{C}ech class of visible centers;
\item every screened asymptotic defect sector is killed by central localization, and only the vacuum sector
survives whenever the visible center is scalar and all nontrivial simples are screened;
\item every Hilbert-screening bridge datum upgrades the visible-center collapse to a concrete spectral-gap criterion,
and a domination estimate for the untwisted Hamiltonian implies a positive mass gap.
\end{enumerate}
\end{theorem}

\begin{proof}
Item (i) is the instanton-completed and screened tangent-to-center package.
Item (ii) is the screened higher-body \v{C}ech duality theorem.
Item (iii) is the central localization criterion for confinement.
Item (iv) is the screening Hodge theorem together with the conditional mass-gap transfer theorem.
\end{proof}

\begin{remark}[Status stratification]
% label removed: rem:status-stratification-platonic-synthesis
The mathematics in this synthesis splits into three strata.
\begin{enumerate}[label=(\alph*)]
\item \emph{Proved here:} boundary BRST recovery, tangent-to-center, central formality,
relative derived centers, higher-body \v{C}ech duality, instanton completion, screening quotients,
visible defect localization, screening Hodge theory, and bounded-error gap stability.
\item \emph{Conditional:} the realization of the dual algebra as the full physical defect algebra,
and the existence of a natural Hilbert-screening bridge for untwisted four-dimensional Yang--Mills.
\item \emph{Conjectural frontier:} a canonical instanton-completed bulk dual, a full untwisting equivalence,
and a proof that the untwisted Hamiltonian is dominated by the screening Laplacian with positive constant.
\end{enumerate}
\end{remark}


% ===== Begin twisted_ym_boundary_package_v2.tex =====
% Insertion note.
% Place this block immediately after Remark~\ref{rem:PW-connection}
% in chapters/connections/holomorphic_topological.tex.
% In the current build, the first new theorem-like item will compile as 31.1.11.
%
% A second optional insertion point for the frontier conjectures is after
% Conjecture~\ref{conj:agt-bar-cobar}; if kept here, the cross-references still compile.

\subsection{Twisted Yang--Mills boundary package}
% label removed: subsec:twisted-ym-boundary-package

\begin{summary}[Programme]
The role of the present subsection is to isolate the mathematically controlled input from
holomorphic-topological Yang--Mills theory that enters the bar-cobar framework.

\begin{enumerate}[label=(\roman*)]
\item \emph{Twisted Yang--Mills input.} Start with a holomorphic-topological twist whose
boundary local operators form a VOA with chiral envelope.
\item \emph{Bar/cobar construction.} Apply the chiral bar and cobar functors of Part~\ref{part:swiss-cheese} to the
boundary envelope.
\item \emph{Boundary VOA.} Use the VOA for explicit OPE computations, but pass to the chiral
envelope for geometric and derived statements.
\item \emph{Line-operator category.} Interpret the Koszul dual as the algebra governing defects
and line operators whenever such an identification is available.
\item \emph{AGT/holography outputs.} Treat AGT and twisted holography as downstream outputs of
the same boundary package, with precise algebraic content and clearly separated conjectural input.
\end{enumerate}
\end{summary}

\begin{definition}[Koszul-admissible twisted Yang--Mills boundary datum]
% label removed: def:twisted-ym-boundary-datum
Let $X$ be a smooth projective curve. A \emph{Koszul-admissible twisted Yang--Mills boundary datum}
on $X$ consists of:
\begin{enumerate}[label=(\roman*)]
\item a holomorphic-topological twist $\mathcal{T}$ of a $4$-dimensional supersymmetric
Yang--Mills theory whose boundary local operators form a VOA in the sense of
Theorem~\ref{thm:paquette-williams-boundaries};
\item a boundary condition $\cB$ with boundary VOA $V_{\cB}$ and chiral envelope $\cA_{\cB}$;
\item the assumption that $\cA_{\cB}$ lies on the Koszul locus of Part~\ref{part:swiss-cheese} and belongs to a class
for which Theorem~\ref{thm:open-string-bar} applies, e.g. generic-level affine
Kac--Moody or principal $\mathcal{W}$-algebra boundary conditions.
\end{enumerate}
\end{definition}

\begin{notation}[Open sector, bulk shadow, and tangent space]
% label removed: not:twisted-ym-boundary-objects
For a Koszul-admissible twisted Yang--Mills boundary datum $(X,\mathcal{T},\cB)$, set
\[
\mathsf{Open}_{\cB} := \barBch(\cA_{\cB}),
\qquad
\mathsf{Bulk}^{\vee}_{\cB} := \Omegach(\cA_{\cB}^!),
\qquad
T^1_X(\cB) := HH^2_{\mathrm{chiral}}(\cA_{\cB}).
\]
We call $\mathsf{Open}_{\cB}$ the \emph{open boundary complex}, $\mathsf{Bulk}^{\vee}_{\cB}$ the
\emph{dual bulk shadow}, and $T^1_X(\cB)$ the \emph{first-order boundary deformation space}.
\end{notation}

\begin{theorem}[Boundary BRST recovery for twisted Yang--Mills data; \ClaimStatusProvedHere]
% label removed: thm:twisted-ym-boundary-brst
Let $(X,\mathcal{T},\cB)$ be a Koszul-admissible twisted Yang--Mills boundary datum.
Then:
\begin{enumerate}[label=(\alph*)]
\item there is a quasi-isomorphism
\[
\mathsf{Open}_{\cB} = \barBch(\cA_{\cB}) \simeq C^{\frac{\infty}{2}+\bullet}(\cA_{\cB}),
\]
so the chiral bar complex computes the open-string BRST complex of the boundary condition;
\item the Verdier dual of the open sector is the bar of the Koszul dual (the homotopy dual algebra):
\[
\mathbb{D}_{\Ran}\, \mathsf{Open}_{\cB} \simeq (\cA_{\cB}^!)_{\infty} = \barBch(\cA_{\cB}^!).
\]
\end{enumerate}
\end{theorem}

\begin{proof}
By Theorem~\ref{thm:paquette-williams-boundaries}, the boundary condition $\cB$ carries a
VOA $V_{\cB}$ with chiral envelope $\cA_{\cB}$. Since $\cA_{\cB}$ is assumed to be
Koszul-admissible, Theorem~\ref{thm:open-string-bar} applies and gives the quasi-isomorphism
$\barBch(\cA_{\cB}) \simeq C^{\frac{\infty}{2}+\bullet}(\cA_{\cB})$. The Verdier-dual statement is
Theorem~\ref*{V1-thm:verdier-bar-cobar} specialized to $\cA_{\cB}$.
\end{proof}

\begin{theorem}[Tangent-to-center principle; \ClaimStatusProvedHere]
% label removed: thm:twisted-ym-tangent-center
For every Koszul-admissible twisted Yang--Mills boundary datum on a smooth projective
curve $X$, there is a canonical isomorphism
\[
T^1_X(\cB)
= HH^2_{\mathrm{chiral}}(\cA_{\cB})
\xrightarrow{\ \sim\ }
Z(\cA_{\cB}^!)^{\vee} \otimes \omega_X.
\]
In particular, the full first-order boundary deformation theory is controlled by the center
of the Koszul dual algebra.
\end{theorem}

\begin{proof}
Apply Corollary~\ref*{V1-cor:def-obs-exchange-genus0} to the Koszul pair
$(\cA_{\cB}, \cA_{\cB}^!)$.
\end{proof}

\begin{corollary}[Center-vanishing rigidity criterion; \ClaimStatusProvedHere]
% label removed: cor:twisted-ym-center-rigidity
If $Z(\cA_{\cB}^!) = 0$, then $HH^2_{\mathrm{chiral}}(\cA_{\cB}) = 0$.
Hence the boundary condition $\cB$ is infinitesimally rigid.
\end{corollary}

\begin{proof}
Immediate from Theorem~\ref{thm:twisted-ym-tangent-center}.
\end{proof}

\begin{corollary}[One-parameter criterion; \ClaimStatusProvedHere]
% label removed: cor:twisted-ym-one-parameter
If $\dim Z(\cA_{\cB}^!) = 1$, then $\dim HH^2_{\mathrm{chiral}}(\cA_{\cB}) = 1$.
Thus the boundary theory has a unique first-order coupling up to scale.
\end{corollary}

\begin{proof}
Again apply Theorem~\ref{thm:twisted-ym-tangent-center}.
\end{proof}

\begin{definition}[Boundary-changing morphism]
% label removed: def:twisted-ym-boundary-morphism
A \emph{boundary-changing morphism}
$f \colon (X,\mathcal{T},\cB) \to (X,\mathcal{T},\cB')$
is a morphism of boundary conditions that induces a morphism of chiral envelopes
\[
f \colon \cA_{\cB} \longrightarrow \cA_{\cB'}.
\]
Let $\mathsf{KBdy}_{\mathrm{YM}}(X)$ denote the resulting category of Koszul-admissible
boundary data on~$X$.
\end{definition}

\begin{theorem}[The tangent functor factors through the dual center; \ClaimStatusProvedHere]
% label removed: thm:twisted-ym-tangent-factors-through-center
The assignment
\[
T^1_X \colon \mathsf{KBdy}_{\mathrm{YM}}(X) \longrightarrow \Vect,
\qquad
\cB \longmapsto HH^2_{\mathrm{chiral}}(\cA_{\cB}),
\]
is naturally isomorphic to the functor
\[
\mathfrak{Z}^{\vee}_{X} \colon \mathsf{KBdy}_{\mathrm{YM}}(X) \longrightarrow \Vect,
\qquad
\cB \longmapsto Z(\cA_{\cB}^!)^{\vee} \otimes \omega_X.
\]
Equivalently, the first-order boundary deformation functor factors through the center of
the Koszul dual.
\end{theorem}

\begin{proof}
Theorem~\ref{thm:twisted-ym-tangent-center} gives the objectwise identification.
For morphisms, Theorem~\ref*{V1-thm:bv-functor} supplies functoriality of the bar construction,
and the construction of Theorem~\ref*{V1-thm:main-koszul-hoch} (Vol~I Theorem~H, at generic
level; the critical level $k = -h^\vee$ is excluded because $\dim \ChirHoch^0$ can be infinite
there, see Vol~I Remark~\ref*{rem:critical-level-lie-vs-chirhoch}) is functorial because it is
built from natural operations: bar/cobar, derived \(\mathrm{Hom}\), and Verdier duality.
Passing to cohomological degree~$2$ produces the required natural transformation. Koszul-admissible
twisted Yang--Mills boundary data are taken at generic level; the critical-level locus is outside
the scope of this functoriality argument.
\end{proof}

\begin{corollary}[Interface invariance of tangent data; \ClaimStatusProvedHere]
% label removed: cor:twisted-ym-interface-invariance
Let $f \colon \cA_{\cB} \to \cA_{\cB'}$ be a boundary-changing morphism.
If the induced map on dual centers
\[
Z(\cA_{\cB'}^!) \xrightarrow{\ \sim\ } Z(\cA_{\cB}^!)
\]
is an isomorphism, then the induced map on first-order deformation spaces
\[
HH^2_{\mathrm{chiral}}(\cA_{\cB}) \xrightarrow{\ \sim\ } HH^2_{\mathrm{chiral}}(\cA_{\cB'})
\]
is an isomorphism as well.
\end{corollary}

\begin{proof}
This is immediate from Theorem~\ref{thm:twisted-ym-tangent-factors-through-center}.
\end{proof}

\begin{corollary}[Comparison without full boundary equivalence; \ClaimStatusProvedHere]
% label removed: cor:twisted-ym-center-comparison
If two Koszul-admissible twisted Yang--Mills boundary conditions $\cB_1$ and $\cB_2$ satisfy
\[
Z(\cA_{\cB_1}^!) \cong Z(\cA_{\cB_2}^!),
\]
then their first-order deformation spaces are canonically isomorphic:
\[
HH^2_{\mathrm{chiral}}(\cA_{\cB_1}) \cong HH^2_{\mathrm{chiral}}(\cA_{\cB_2}).
\]
In particular, first-order boundary moduli are a strictly coarser invariant than the full
boundary VOA.
\end{corollary}

\begin{proof}
Apply Theorem~\ref{thm:twisted-ym-tangent-center} to each boundary condition.
\end{proof}

\begin{remark}[Mathematical content and physical content]
% label removed: rem:twisted-ym-math-vs-physics
Theorems~\ref{thm:twisted-ym-boundary-brst}--\ref{thm:twisted-ym-tangent-factors-through-center}
and Corollaries~\ref{cor:twisted-ym-center-rigidity}--\ref{cor:twisted-ym-center-comparison}
are entirely internal to the present manuscript once the boundary algebra is known to lie on
the Koszul-admissible locus. The genuinely physical input is the production of such a
boundary algebra from a holomorphic-topological Yang--Mills theory.
\end{remark}
% New frontier mathematics package: central formality and mixed couplings.
% Suggested insertion point: immediately after the subsection
% "Twisted Yang--Mills boundary package" in Chapter 31.

\subsection{Central formality and mixed-coupling theorems}
% label removed: subsec:twisted-ym-central-formality

\begin{definition}[Boundary deformation dg Lie algebra]
% label removed: def:twisted-ym-def-dglie
Let $(X,\mathcal{T},\cB)$ be a Koszul-admissible twisted Yang--Mills boundary datum, with chiral envelope $\cA_{\cB}$. Define the \emph{boundary deformation dg Lie algebra}
\[
\mathfrak{g}^{\mathrm{def}}_{\cB}
:=
\tau_{\ge 1} RHH_{\mathrm{ch}}(\cA_{\cB})[1],
\]
where $RHH_{\mathrm{ch}}(\cA_{\cB})$ carries the standard shifted dg Lie structure of the chiral Hochschild complex. Its Maurer--Cartan functor governs formal boundary deformations of $\cA_{\cB}$.
\end{definition}

\begin{theorem}[Dual unobstructedness criterion; \ClaimStatusProvedHere]
% label removed: thm:twisted-ym-dual-unobstructedness
Let $(X,\mathcal{T},\cB)$ be a Koszul-admissible twisted Yang--Mills boundary datum on a smooth projective curve. If
\[
HH^{-1}_{\mathrm{chiral}}(\cA_{\cB}^{!}) = 0,
\]
then the primary obstruction space of $\cA_{\cB}$ vanishes:
\[
HH^{3}_{\mathrm{chiral}}(\cA_{\cB}) = 0.
\]
Hence every first-order deformation of $\cA_{\cB}$ is unobstructed to second order.
\end{theorem}

\begin{proof}
By the deformation-obstruction dictionary, primary obstructions to extending a first-order deformation lie in $HH^{3}_{\mathrm{chiral}}(\cA_{\cB})$. Apply Theorem~\ref*{V1-thm:main-koszul-hoch} (Vol~I Theorem~H, at generic level; the critical level $k = -h^\vee$ is excluded because $\dim \ChirHoch^0$ can be infinite there, see Vol~I Remark~\ref*{rem:critical-level-lie-vs-chirhoch}) with $n=3$ to obtain
\[
HH^{3}_{\mathrm{chiral}}(\cA_{\cB})
\cong
HH^{-1}_{\mathrm{chiral}}(\cA_{\cB}^{!})^{\vee}\otimes \omega_X.
\]
The assumed vanishing of $HH^{-1}_{\mathrm{chiral}}(\cA_{\cB}^{!})$ forces $HH^{3}_{\mathrm{chiral}}(\cA_{\cB})=0$, so the primary obstruction vanishes. The final sentence is the standard interpretation of $HH^3$ in deformation theory. The generic-level restriction propagates: the dual unobstructedness criterion is asserted for Koszul-admissible boundary data at generic level.
\end{proof}

\begin{theorem}[Central formality theorem; \ClaimStatusProvedHere]
% label removed: thm:twisted-ym-central-formality
Let $(X,\mathcal{T},\cB)$ be a Koszul-admissible twisted Yang--Mills boundary datum. Assume
\[
HH^{1}_{\mathrm{chiral}}(\cA_{\cB}) = 0
\qquad\text{and}\qquad
HH^{-1}_{\mathrm{chiral}}(\cA_{\cB}^{!}) = 0.
\]
Then the formal deformation problem of $\cA_{\cB}$ is equivalent to the formal affine space on the vector space
\[
T^1_X(\cB)
\cong
Z(\cA_{\cB}^{!})^{\vee}\otimes \omega_X.
\]
Equivalently, every first-order boundary deformation extends uniquely to all orders up to unique equivalence, and the completed local ring of the deformation space is
\[
\widehat{\mathrm{Sym}}\bigl(Z(\cA_{\cB}^{!})\otimes \omega_X^{-1}\bigr).
\]
\end{theorem}

\begin{proof}
By Theorem~\ref{thm:MC-deformations}, the formal deformations of $\cA_{\cB}$ are governed by Maurer--Cartan elements of the dg Lie algebra $\mathfrak{g}^{\mathrm{def}}_{\cB}$. Its cohomology satisfies
\[
H^0(\mathfrak{g}^{\mathrm{def}}_{\cB}) = HH^{1}_{\mathrm{chiral}}(\cA_{\cB}) = 0,
\qquad
H^1(\mathfrak{g}^{\mathrm{def}}_{\cB}) = HH^{2}_{\mathrm{chiral}}(\cA_{\cB}),
\]
and, by Theorem~\ref{thm:twisted-ym-dual-unobstructedness},
\[
H^2(\mathfrak{g}^{\mathrm{def}}_{\cB}) = HH^{3}_{\mathrm{chiral}}(\cA_{\cB}) = 0.
\]
Choose a minimal $L_\infty$ model on $H^\bullet(\mathfrak{g}^{\mathrm{def}}_{\cB})$; existence follows from the homotopy transfer theorem applied to the Lie operad, and uniqueness up to $L_\infty$ isomorphism is standard. Since the only nonzero cohomology relevant to the Maurer--Cartan problem is in degree $1$, every higher bracket
\[
\ell_n \colon (H^1)^{\otimes n} \longrightarrow H^{2}
\]
has zero target, because $\ell_n$ has cohomological degree $2-n$. Thus all $\ell_n$ vanish identically and the minimal model is abelian.

The Maurer--Cartan equation in an abelian dg Lie algebra is empty, and the gauge action is trivial because $H^0=0$. Therefore the formal moduli problem is the formal affine space on $H^1(\mathfrak{g}^{\mathrm{def}}_{\cB}) = HH^2_{\mathrm{chiral}}(\cA_{\cB})$. Finally, Theorem~\ref{thm:twisted-ym-tangent-center} identifies this tangent space with $Z(\cA_{\cB}^{!})^{\vee}\otimes\omega_X$, yielding the completed symmetric algebra stated above.
\end{proof}

\begin{remark}[Interpretation]
% label removed: rem:twisted-ym-central-formality-meaning
The theorem says that under vanishing of infinitesimal automorphisms and dual negative-degree Hochschild classes, the entire formal boundary moduli problem linearizes. The Koszul dual center is not merely the tangent space: it gives actual formal coordinates on the deformation space.
\end{remark}

\begin{definition}[Reduced dual center and reduced tangent space]
% label removed: def:twisted-ym-reduced-center
Assume that the boundary algebra is augmented, so that the vacuum defines a distinguished scalar subspace $\mathbb{C}\mathbf{1} \subset Z(\cA_{\cB}^{!})$. Define the \emph{reduced dual center}
\[
\widetilde Z(\cA_{\cB}^{!}) := Z(\cA_{\cB}^{!})/\mathbb{C}\mathbf{1}
\]
and the \emph{reduced tangent space}
\[
\widetilde T^1_X(\cB) := \widetilde Z(\cA_{\cB}^{!})^{\vee}\otimes \omega_X.
\]
\end{definition}

\begin{lemma}[Center of a chiral tensor product; \ClaimStatusProvedHere]
% label removed: lem:twisted-ym-center-tensor
Let $\cB_1,\cB_2$ be quadratic Koszul-admissible boundary data with chiral envelopes $\cA_1,\cA_2$, and assume that the chiral tensor product $\cA_1\otimes^{\mathrm{ch}}\cA_2$ has trivial mixed OPE \/in the sense of Proposition~\ref*{V1-prop:koszul-dual-tensor-product}. Then
\[
Z\bigl(\cA_1\otimes^{\mathrm{ch}}\cA_2\bigr)
\cong
Z(\cA_1)\otimes Z(\cA_2).
\]
\end{lemma}

\begin{proof}
An element of $\cA_1\otimes^{\mathrm{ch}}\cA_2$ is central iff it has regular OPE with every generator of the form $a\otimes \mathbf{1}$ and $\mathbf{1}\otimes b$. Because the mixed OPE is trivial, regularity against all $a\otimes \mathbf{1}$ is equivalent to lying in $Z(\cA_1)\otimes \cA_2$, while regularity against all $\mathbf{1}\otimes b$ is equivalent to lying in $\cA_1\otimes Z(\cA_2)$. Hence
\[
Z(\cA_1\otimes^{\mathrm{ch}}\cA_2)
=
\bigl(Z(\cA_1)\otimes \cA_2\bigr)
\cap
\bigl(\cA_1\otimes Z(\cA_2)\bigr).
\]
Choose vector space splittings $\cA_i = Z(\cA_i)\oplus U_i$. Then
\[
Z(\cA_1)\otimes \cA_2
=
Z(\cA_1)\otimes Z(\cA_2)
\oplus
Z(\cA_1)\otimes U_2,
\]
while
\[
\cA_1\otimes Z(\cA_2)
=
Z(\cA_1)\otimes Z(\cA_2)
\oplus
U_1\otimes Z(\cA_2).
\]
Their intersection is exactly $Z(\cA_1)\otimes Z(\cA_2)$.
\end{proof}

\begin{theorem}[Binary mixed-coupling theorem; \ClaimStatusProvedHere]
% label removed: thm:twisted-ym-binary-mixed-couplings
Let $\cB_1,\cB_2$ be quadratic Koszul-admissible twisted Yang--Mills boundary data on the same smooth projective curve $X$, with chiral envelopes $\cA_1,\cA_2$. Then the composite boundary datum with envelope
\[
\cA_{12}:=\cA_1\otimes^{\mathrm{ch}}\cA_2
\]
satisfies
\[
T^1_X(\cB_{12})
\cong
\bigl(Z(\cA_1^{!})\otimes Z(\cA_2^{!})\bigr)^{\vee}\otimes \omega_X.
\]
If the centers are finite-dimensional and split as
\[
Z(\cA_i^{!}) \cong \mathbb{C}\mathbf{1}\oplus \widetilde Z_i,
\]
then there is a canonical decomposition
\[
\widetilde T^1_X(\cB_{12})
\cong
\widetilde T^1_X(\cB_1)
\oplus
\widetilde T^1_X(\cB_2)
\oplus
(\widetilde Z_1\otimes \widetilde Z_2)^{\vee}\otimes \omega_X.
\]
The last summand is a genuine mixed first-order coupling, invisible in either factor separately.
\end{theorem}

\begin{proof}
By Theorem~\ref{thm:twisted-ym-tangent-center},
\[
T^1_X(\cB_{12})
\cong
Z(\cA_{12}^{!})^{\vee}\otimes \omega_X.
\]
Since the boundary data are quadratic, Proposition~\ref*{V1-prop:koszul-dual-tensor-product} gives
\[
\cA_{12}^{!}
\cong
\cA_1^{!}\otimes^{\mathrm{ch}}\cA_2^{!}.
\]
Applying Lemma~\ref{lem:twisted-ym-center-tensor} to the dual side yields
\[
Z(\cA_{12}^{!})
\cong
Z(\cA_1^{!})\otimes Z(\cA_2^{!}),
\]
which proves the first statement.

For the decomposition, write each center as $\mathbb{C}\mathbf{1}\oplus \widetilde Z_i$. Then
\[
Z(\cA_1^{!})\otimes Z(\cA_2^{!})
\cong
\mathbb{C}\mathbf{1}
\oplus
\widetilde Z_1
\oplus
\widetilde Z_2
\oplus
(\widetilde Z_1\otimes \widetilde Z_2).
\]
Dualizing and tensoring with $\omega_X$ gives the asserted decomposition of reduced tangent spaces.
\end{proof}

\begin{corollary}[Multi-body coupling expansion; \ClaimStatusProvedHere]
% label removed: cor:twisted-ym-multibody-couplings
Let $\cB_1,\dots,\cB_m$ be quadratic Koszul-admissible twisted Yang--Mills boundary data on $X$, and suppose each dual center splits as
\[
Z(\cA_i^{!}) \cong \mathbb{C}\mathbf{1}\oplus \widetilde Z_i.
\]
Then the reduced first-order deformation space of the composite boundary datum
\[
\cB_{1\cdots m}
:=
\cB_1\otimes^{\mathrm{ch}}\cdots\otimes^{\mathrm{ch}}\cB_m
\]
admits the canonical expansion
\[
\widetilde T^1_X(\cB_{1\cdots m})
\cong
\bigoplus_{\emptyset \neq S \subseteq \{1,\dots,m\}}
\left(\bigotimes_{i\in S} \widetilde Z_i\right)^{\vee}\otimes \omega_X.
\]
In particular, subsets $S$ of size $r$ parameterize genuine $r$-body first-order couplings.
\end{corollary}

\begin{proof}
Iterate Theorem~\ref{thm:twisted-ym-binary-mixed-couplings}. Inductively one obtains
\[
Z(\cA_{1\cdots m}^{!})
\cong
\bigotimes_{i=1}^m Z(\cA_i^{!})
\cong
\bigotimes_{i=1}^m (\mathbb{C}\mathbf{1}\oplus \widetilde Z_i).
\]
Expanding the tensor product yields a direct sum over subsets $S\subseteq \{1,\dots,m\}$. Removing the empty subset gives the reduced statement, and dualizing/tensoring with $\omega_X$ proves the formula.
\end{proof}

\begin{corollary}[Mixed rigidity criterion; \ClaimStatusProvedHere]
% label removed: cor:twisted-ym-mixed-rigidity
Under the hypotheses of Corollary~\ref{cor:twisted-ym-multibody-couplings}:
\begin{enumerate}[label=(\roman*)]
\item if every $\widetilde Z_i=0$, then every composite boundary condition is reduced-infinitesimally rigid;
\item if exactly one $\widetilde Z_i$ is nonzero, then there are no genuine mixed couplings;
\item a nonzero tensor product $\widetilde Z_{i_1}\otimes\cdots\otimes\widetilde Z_{i_r}$ is both necessary and sufficient for the existence of an $r$-body first-order coupling.
\end{enumerate}
\end{corollary}

\begin{proof}
This is immediate from Corollary~\ref{cor:twisted-ym-multibody-couplings} by inspection of the subset expansion.
\end{proof}

\begin{remark}[Why this is genuinely new]
% label removed: rem:twisted-ym-why-new
Theorems~\ref{thm:twisted-ym-dual-unobstructedness}--\ref{thm:twisted-ym-binary-mixed-couplings} do not merely repackage the tangent-to-center principle. They add three new pieces of structure:
\begin{enumerate}[label=(\roman*)]
\item a \emph{dual-side obstruction criterion} in negative Hochschild degree;
\item a \emph{formality theorem} identifying explicit formal coordinates on the boundary deformation space;
\item a \emph{multi-body coupling law} showing that composite boundary theories acquire genuinely mixed couplings controlled by tensor products of reduced dual centers.
\end{enumerate}
These statements push the present framework from a static duality theorem to an actual local geometry of boundary moduli.
\end{remark}
% ===== End {fn} =====

% [Duplicate block removed -- content consolidated above at line 279.]


% ===== Begin w_sector_relative_frontier_package.tex =====
% Frontier package: relative Hochschild duality and mixed Kodaira--Spencer class
% Suggested insertion point: after the subsection on derived mixed couplings for non-quadratic W-boundaries.

\subsection{Relative Hochschild duality and the mixed Kodaira--Spencer class}
% label removed: subsec:relative-hochschild-mixed-ks

Fix a smooth projective curve $X$ and a weakly composable pair of transferred $W$-dual data
\[
M_1 \xrightarrow{\iota_1} M_{12} \xleftarrow{\iota_2} M_2,
\]
where each $M_i$ is a complete filtered minimal $A_\infty$-algebra quasi-isomorphic to the transferred dual of a principal $W$-boundary condition.

\begin{definition}[Relative Hochschild cochains]
% label removed: def:relative-hochschild-cochains
Define the \emph{relative Hochschild cochain complex}
\[
CC^\bullet_{\mathrm{rel}}(M_{12};M_1,M_2)
:=
\operatorname{hocofib}\!\Big(
CC^\bullet(M_1)\oplus CC^\bullet(M_2)
\xrightarrow{\;\iota_{1*}+\iota_{2*}\;}
CC^\bullet(M_{12})
\Big).
\]
Its cohomology groups are denoted
\[
HH^\bullet_{\mathrm{rel}}(M_{12};M_1,M_2)
:= H^\bullet\big(CC^\bullet_{\mathrm{rel}}(M_{12};M_1,M_2)\big).
\]
The degree-zero group is called the \emph{relative derived center}.
\end{definition}

\begin{definition}[Mixed deformation $L_\infty$-algebra]
% label removed: def:mixed-deformation-linfty-relative
Let
\[
\mathfrak g_i:=\tau_{\ge 1}CC^\bullet(M_i)[1],
\qquad
\mathfrak g_{12}:=\tau_{\ge 1}CC^\bullet(M_{12})[1],
\]
with the standard filtered $L_\infty$ structures induced by the Gerstenhaber bracket. The restriction maps induced by $\iota_1,\iota_2$ define an $L_\infty$ morphism
\[
r=(r_1,r_2):\mathfrak g_{12}\to \mathfrak g_1\oplus \mathfrak g_2.
\]
Define the \emph{mixed deformation $L_\infty$-algebra}
\[
\mathfrak g_{\mathrm{mix}}:=\operatorname{hofib}(r).
\]
\end{definition}

\begin{theorem}[Relative duality theorem; \ClaimStatusProvedHere]
% label removed: thm:relative-duality-theorem
Assume the degreewise complementarity isomorphisms
\[
H^k(\mathfrak g_i)
\cong
HH^{1-k}(M_i)^\vee\otimes \omega_X,
\qquad
H^k(\mathfrak g_{12})
\cong
HH^{1-k}(M_{12})^\vee\otimes \omega_X
\]
for all $k$ under consideration, and assume the relevant cohomology groups are finite-dimensional. Then for every integer $k$ there is a natural isomorphism
\[
H^k(\mathfrak g_{\mathrm{mix}})
\cong
HH^{1-k}_{\mathrm{rel}}(M_{12};M_1,M_2)^\vee\otimes \omega_X.
\]
In particular,
\[
H^1(\mathfrak g_{\mathrm{mix}})
\cong
HH^{0}_{\mathrm{rel}}(M_{12};M_1,M_2)^\vee\otimes \omega_X,
\]
and
\[
H^2(\mathfrak g_{\mathrm{mix}})
\cong
HH^{-1}_{\mathrm{rel}}(M_{12};M_1,M_2)^\vee\otimes \omega_X.
\]
\end{theorem}

\begin{proof}
The homotopy fiber defining $\mathfrak g_{\mathrm{mix}}$ gives a distinguished triangle
\[
\mathfrak g_{\mathrm{mix}} \to \mathfrak g_{12} \xrightarrow{r} \mathfrak g_1\oplus \mathfrak g_2 \to \mathfrak g_{\mathrm{mix}}[1],
\]
hence a long exact sequence in cohomology. The homotopy cofiber defining $CC^\bullet_{\mathrm{rel}}$ gives a distinguished triangle
\[
CC^\bullet(M_1)\oplus CC^\bullet(M_2)
\to CC^\bullet(M_{12})
\to CC^\bullet_{\mathrm{rel}}(M_{12};M_1,M_2)
\to
\big(CC^\bullet(M_1)\oplus CC^\bullet(M_2)\big)[1],
\]
and therefore a long exact sequence in Hochschild cohomology.

By the assumed complementarity isomorphisms, the cohomology groups of the first triangle are dual, up to the twist by $\omega_X$, to the cohomology groups of the second triangle with the degree shift $k\mapsto 1-k$. Naturality of the restriction maps identifies the boundary morphisms in the two long exact sequences. Dualizing the second long exact sequence and comparing term-by-term with the first, the five-lemma yields the stated isomorphisms for all $k$.
\end{proof}

\begin{corollary}[Relative tangent-to-center principle; \ClaimStatusProvedHere]
% label removed: cor:relative-tangent-center-principle
The space of genuine mixed first-order couplings satisfies
\[
T^1_{X,\mathrm{mix}}(\cB_{12};\cB_1,\cB_2)
:= H^1(\mathfrak g_{\mathrm{mix}})
\cong
HH^{0}_{\mathrm{rel}}(M_{12};M_1,M_2)^\vee\otimes \omega_X.
\]
\end{corollary}

\begin{corollary}[Relative dual unobstructedness; \ClaimStatusProvedHere]
% label removed: cor:relative-dual-unobstructedness
If
\[
HH^{-1}_{\mathrm{rel}}(M_{12};M_1,M_2)=0,
\]
then every genuine mixed first-order coupling extends to second order.
\end{corollary}

\begin{proof}
By Theorem~\ref{thm:relative-duality-theorem}, the primary mixed obstruction space is
\[
H^2(\mathfrak g_{\mathrm{mix}})
\cong
HH^{-1}_{\mathrm{rel}}(M_{12};M_1,M_2)^\vee\otimes \omega_X.
\]
Thus the assumed vanishing forces the obstruction space to vanish.
\end{proof}

\begin{theorem}[Relative central formality; \ClaimStatusProvedHere]
% label removed: thm:relative-central-formality
Assume
\[
HH^{1}_{\mathrm{rel}}(M_{12};M_1,M_2)=0
\qquad\text{and}\qquad
HH^{-1}_{\mathrm{rel}}(M_{12};M_1,M_2)=0.
\]
Then the formal moduli problem of genuine mixed deformations is equivalent to the formal affine space on
\[
HH^{0}_{\mathrm{rel}}(M_{12};M_1,M_2)^\vee\otimes \omega_X.
\]
Equivalently, every mixed first-order coupling extends uniquely to all orders up to unique equivalence.
\end{theorem}

\begin{proof}
By Theorem~\ref{thm:relative-duality-theorem},
\[
H^0(\mathfrak g_{\mathrm{mix}})
\cong HH^{1}_{\mathrm{rel}}(M_{12};M_1,M_2)^\vee\otimes \omega_X=0,
\]
\[
H^1(\mathfrak g_{\mathrm{mix}})
\cong HH^{0}_{\mathrm{rel}}(M_{12};M_1,M_2)^\vee\otimes \omega_X,
\]
and
\[
H^2(\mathfrak g_{\mathrm{mix}})
\cong HH^{-1}_{\mathrm{rel}}(M_{12};M_1,M_2)^\vee\otimes \omega_X=0.
\]
Transfer the $L_\infty$ structure on $\mathfrak g_{\mathrm{mix}}$ to a minimal model on its cohomology. Degree counting forces every higher bracket on degree-one classes to land in degree two, which vanishes. Thus the minimal model is abelian on the deformation sector, the Maurer--Cartan equation is empty, and the gauge action is trivial because $H^0=0$.
\end{proof}

\begin{definition}[Mixed Kodaira--Spencer class]
% label removed: def:mixed-kodaira-spencer-class
Suppose $M_{12}$ carries a complete descending filtration by mixed weight, and write its $A_\infty$ coderivation as
\[
m = m^{(0)} + m^{(1)} + m^{(2)} + \cdots,
\]
where $m^{(r)}$ raises mixed weight by $r$, and $m^{(0)}$ is the decoupled part inherited from $M_1$ and $M_2$.

The class of $m^{(1)}$ in the degree-one relative Hochschild cohomology,
\[
\kappa_{\mathrm{mix}}(M_{12};M_1,M_2)
:= [m^{(1)}]
\in HH^{1}_{\mathrm{rel}}(M_{12};M_1,M_2),
\]
is called the \emph{mixed Kodaira--Spencer class}.
\end{definition}

\begin{lemma}[Well-definedness and invariance; \ClaimStatusProvedHere]
% label removed: lem:mixed-ks-well-defined
The mixed Kodaira--Spencer class is independent of the choice of filtered minimal model, up to the canonical identification of relative Hochschild cohomology under filtered $A_\infty$ quasi-isomorphism.
\end{lemma}

\begin{proof}
A filtered $A_\infty$ quasi-isomorphism between two transferred models intertwines their coderivations up to gauge equivalence in the Hochschild dg Lie algebra. The weight-one component of a gauge-equivalent coderivation changes by a relative Hochschild coboundary. Hence the cohomology class of $m^{(1)}$ is invariant.
\end{proof}

\begin{theorem}[First-correction theorem; \ClaimStatusProvedHere]
% label removed: thm:first-correction-theorem
Equip $CC^\bullet_{\mathrm{rel}}(M_{12};M_1,M_2)$ with the filtration induced by mixed weight. Then the associated spectral sequence has
\[
E_1^0 \cong HH^0\bigl(\operatorname{gr} CC^\bullet_{\mathrm{rel}}\bigr),
\]
and the first differential is induced by the adjoint action of the mixed Kodaira--Spencer class:
\[
d_1[\phi] = [\kappa_{\mathrm{mix}},\phi].
\]
In particular, on the quadratic-limit term
\[
E_1^0 \cong \widetilde Z_\infty(M_1)\widehat\otimes \widetilde Z_\infty(M_2),
\]
the class of a naive quadratic mixed coupling survives to the next page if and only if it lies in the kernel of
\[
[\kappa_{\mathrm{mix}},-] :
\widetilde Z_\infty(M_1)\widehat\otimes \widetilde Z_\infty(M_2)
\longrightarrow E_1^1.
\]
\end{theorem}

\begin{proof}
On Hochschild cochains, the differential is the Gerstenhaber commutator with the full coderivation:
\[
d = [m,-] = [m^{(0)},-] + [m^{(1)},-] + [m^{(2)},-] + \cdots.
\]
The zeroth differential on the associated graded is therefore $d_0=[m^{(0)},-]$, and the first derived differential is induced by the next term $[m^{(1)},-]$. By Definition~\ref{def:mixed-kodaira-spencer-class}, the cohomology class of $m^{(1)}$ is precisely $\kappa_{\mathrm{mix}}$, so on the $E_1$-page the first differential is the adjoint action of $\kappa_{\mathrm{mix}}$.
\end{proof}

\begin{corollary}[Persistence criterion for the quadratic law; \ClaimStatusProvedHere]
% label removed: cor:persistence-criterion-quadratic-law
The quadratic mixed-coupling law survives to first non-quadratic order if and only if
\[
\kappa_{\mathrm{mix}}(M_{12};M_1,M_2)=0.
\]
More generally, if the mixed coderivation components $m^{(1)},\dots,m^{(r)}$ vanish in relative Hochschild cohomology, then the mixed-weight spectral sequence degenerates through $E_{r+1}$.
\end{corollary}

\begin{proof}
The first statement is immediate from Theorem~\ref{thm:first-correction-theorem}. The second follows from the same filtered-complex argument: if all differential pieces up to mixed weight $r$ are cohomologically trivial, then all differentials $d_1,\dots,d_r$ vanish.
\end{proof}

\begin{theorem}[Universal interaction brackets on mixed couplings; \ClaimStatusProvedHere]
% label removed: thm:universal-interaction-brackets
Let
\[
\mathbb T_{\mathrm{mix}}:=H^1(\mathfrak g_{\mathrm{mix}}),
\qquad
\mathrm{Obs}_{\mathrm{mix}}:=H^2(\mathfrak g_{\mathrm{mix}}).
\]
Homotopy transfer endows $H^\bullet(\mathfrak g_{\mathrm{mix}})$ with a minimal $L_\infty$ structure whose restriction to degree one defines multilinear maps
\[
\lambda_n : \operatorname{Sym}^n(\mathbb T_{\mathrm{mix}}) \longrightarrow \mathrm{Obs}_{\mathrm{mix}},
\qquad n\ge 2.
\]
A formal mixed coupling $\alpha\in \mathbb T_{\mathrm{mix}}\widehat\otimes \mathfrak m$ extends to all orders if and only if
\[
\sum_{n\ge 2}\frac{1}{n!}\lambda_n(\alpha,\dots,\alpha)=0.
\]
The first nonzero $\lambda_n$ is a derived invariant of the filtered $A_\infty$ quasi-isomorphism class of the cospan $(M_1\to M_{12}\leftarrow M_2)$.
\end{theorem}

\begin{proof}
This is the standard Maurer--Cartan description for the minimal $L_\infty$ model of the controlling deformation algebra $\mathfrak g_{\mathrm{mix}}$. Since every input has degree $1$ and $L_\infty$ brackets have degree $2-n$, the output lies in degree $2$, i.e. in $\mathrm{Obs}_{\mathrm{mix}}$. Invariance under filtered quasi-isomorphism follows from invariance of the minimal model up to $L_\infty$ isomorphism.
\end{proof}
% ===== End {fn} =====


% ===== Begin twisted_ym_w_boundary_simplex.tex =====
% New frontier package: higher-body Cech theory for non-quadratic W-boundaries.
% This block is self-contained and assumes only the derived tangent-to-center input
% for transferred W-duals.

\subsection{Higher-body \v{C}ech theory for non-quadratic $\mathcal W$-boundaries}
% label removed: subsec:w-boundary-cech-theory

\begin{convention}[Derived tangent-to-center input]
% label removed: conv:derived-tangent-center-input
Throughout this subsection, let $X$ be a smooth projective curve and let
\[
\mathfrak N=\{M_S\}_{\emptyset\neq S\subseteq I},\qquad I=\{1,\dots,m\},
\]
be a finite collection of transferred $\mathcal W$-dual data indexed by nonempty subsets of $I$.
For each $S\subseteq I$, write $\cB_S$ for the corresponding boundary condition and assume there is a canonical identification
\[
T^1_X(\cB_S)
\xrightarrow{\sim}
HH^0(M_S)^\vee\otimes \omega_X.
\]
Set
\[
Z_S:=HH^0(M_S),\qquad T_S:=T^1_X(\cB_S).
\]
\end{convention}

\begin{definition}[Boundary simplex]
% label removed: def:w-boundary-simplex
A \emph{boundary simplex of order $m$} consists of transferred $\mathcal W$-dual data
\[
\mathfrak N=\{M_S\}_{\emptyset\neq S\subseteq I}
\]
with filtered $A_\infty$-morphisms
\[
\rho_{S,T}\colon M_T\longrightarrow M_S
\qquad (\emptyset\neq T\subseteq S\subseteq I)
\]
satisfying $\rho_{S,S}=\id$ and $\rho_{U,S}\circ \rho_{S,T}=\rho_{U,T}$ whenever
$T\subseteq S\subseteq U$.
\end{definition}

\begin{definition}[Derived-center \v{C}ech complex]
% label removed: def:w-derived-center-cech
Fix the natural order on $I$. For a boundary simplex $\mathfrak N$, define the cochain complex
\[
C_Z^p(\mathfrak N):=\bigoplus_{|S|=p+1} Z_S,
\qquad 0\le p\le m-1,
\]
with differential $\delta\colon C_Z^p\to C_Z^{p+1}$ given on the component indexed by
$S=\{s_0<\cdots<s_{p+1}\}$ by
\[
(\delta z)_S
:=
\sum_{j=0}^{p+1}(-1)^j (\rho_{S,S\setminus\{s_j\}})_*(z_{S\setminus\{s_j\}}).
\]
We call $C_Z^\bullet(\mathfrak N)$ the \emph{derived-center \v{C}ech complex} of the simplex.
\end{definition}

\begin{definition}[Tangent restriction complex and pure $m$-body couplings]
% label removed: def:w-tangent-restriction-complex
Define the cochain complex
\[
C_T^p(\mathfrak N):=\bigoplus_{|S|=m-p} T_S,
\qquad 0\le p\le m-1,
\]
with differential $\partial\colon C_T^p\to C_T^{p+1}$ given on the component indexed by
$T=\{t_0<\cdots<t_{m-p-2}\}$ by
\[
(\partial\lambda)_T
:=
\sum_{i\notin T}(-1)^{\operatorname{pos}(i,T\cup\{i\})-1}
\operatorname{res}_{T\cup\{i\},T}(\lambda_{T\cup\{i\}}),
\]
where $\operatorname{res}_{S,T}\colon T_S\to T_T$ is the restriction map dual to
$(\rho_{S,T})_*\colon Z_T\to Z_S$ under Convention~\ref{conv:derived-tangent-center-input}.

The space of \emph{pure $m$-body first-order couplings} is
\[
T^{1,[m]}_{X,\mathrm{pure}}(\mathfrak N):=H^0\bigl(C_T^\bullet(\mathfrak N)\bigr).
\]
\end{definition}

\begin{lemma}[Homotopy invariance; \ClaimStatusProvedHere]
% label removed: lem:w-boundary-simplex-homotopy-invariance
Let $F\colon \mathfrak N\to \mathfrak N'$ be a morphism of boundary simplices which is an
$A_\infty$-quasi-isomorphism on every face and commutes with all face maps. Then the induced map
\[
C_Z^\bullet(\mathfrak N)\longrightarrow C_Z^\bullet(\mathfrak N')
\]
is an isomorphism of cochain complexes. In particular,
\[
H^p\bigl(C_Z^\bullet(\mathfrak N)\bigr)
\cong
H^p\bigl(C_Z^\bullet(\mathfrak N')\bigr)
\]
for all $p$, and the pure coupling spaces are canonically identified.
\end{lemma}

\begin{proof}
Degree-zero Hochschild cohomology is invariant under $A_\infty$-quasi-isomorphism. Hence
$Z_S\cong Z'_S$ for every face $S$, compatibly with the induced maps on $HH^0$. Taking direct sums over all faces of a fixed cardinality yields an isomorphism in each cochain degree, and the compatibility with face maps makes this an isomorphism of complexes. The final statement follows from Theorem~\ref{thm:w-boundary-cech-duality} below.
\end{proof}

\begin{theorem}[Higher-body \v{C}ech duality; \ClaimStatusProvedHere]
% label removed: thm:w-boundary-cech-duality
For every boundary simplex $\mathfrak N$ of order $m$, there is a canonical isomorphism of cochain complexes
\[
C_T^p(\mathfrak N)
\xrightarrow{\sim}
\Bigl(C_Z^{m-1-p}(\mathfrak N)\Bigr)^\vee\otimes \omega_X,
\qquad 0\le p\le m-1,
\]
under which the differential $\partial$ is the transpose of $\delta$.
Consequently, for every $0\le q\le m-1$,
\[
H^q\bigl(C_T^\bullet(\mathfrak N)\bigr)
\cong
H^{m-1-q}\bigl(C_Z^\bullet(\mathfrak N)\bigr)^\vee\otimes \omega_X.
\]
In particular,
\[
T^{1,[m]}_{X,\mathrm{pure}}(\mathfrak N)
\cong
H^{m-1}\bigl(C_Z^\bullet(\mathfrak N)\bigr)^\vee\otimes \omega_X.
\]
\end{theorem}

\begin{proof}
By Convention~\ref{conv:derived-tangent-center-input}, each summand $T_S$ is canonically
identified with $Z_S^\vee\otimes \omega_X$. Summing over all subsets $S$ of cardinality $m-p$
gives the objectwise isomorphism
\[
C_T^p(\mathfrak N)
\cong
\Bigl(C_Z^{m-1-p}(\mathfrak N)\Bigr)^\vee\otimes \omega_X.
\]
The sign conventions in Definitions~\ref{def:w-derived-center-cech} and
\ref{def:w-tangent-restriction-complex} were chosen so that the restriction map
$\operatorname{res}_{S,T}$ is exactly the dual transpose of $(\rho_{S,T})_*$. Therefore the alternating
sum defining $\partial$ is the transpose of the alternating sum defining $\delta$.
Taking cohomology yields the stated duality.
\end{proof}

\begin{remark}[Interpretation]
% label removed: rem:w-boundary-cech-interpretation
The groups $H^q(C_T^\bullet(\mathfrak N))$ form a complete tower of first-order descent obstructions.
A degree-$q$ cochain assigns couplings to all faces of codimension $q$; the cocycle condition means
that these assignments are compatible on codimension-$(q+1)$ faces. The cohomology class measures
exactly the failure of such compatible face data to come from a coupling one dimension higher.
Theorem~\ref{thm:w-boundary-cech-duality} identifies this tower with the dual \v{C}ech cohomology of
facewise derived centers.
\end{remark}

\begin{definition}[Decoupled simplex]
% label removed: def:w-decoupled-simplex
A boundary simplex $\mathfrak N$ is called \emph{decoupled} if there exist complete vector spaces
$\widetilde Z_1,\dots,\widetilde Z_m$ such that, for every nonempty subset $S\subseteq I$,
\[
Z_S\cong \widehat\bigotimes_{i\in S}(\Bbbk\mathbf{1}\oplus \widetilde Z_i),
\]
and under these identifications the maps $(\rho_{S,T})_*$ are obtained by adjoining the unit in every
factor belonging to $S\setminus T$.
\end{definition}

\begin{definition}[Augmented center complex]
% label removed: def:w-augmented-center-complex
For a boundary simplex $\mathfrak N$ of order $m$, define the augmented cochain complex
$\widetilde C_Z^\bullet(\mathfrak N)$ by
\[
\widetilde C_Z^{-1}(\mathfrak N):=\Bbbk,
\qquad
\widetilde C_Z^p(\mathfrak N):=C_Z^p(\mathfrak N)\quad (0\le p\le m-1),
\]
with the augmentation map $\Bbbk\to \bigoplus_i Z_{\{i\}}$ given by the direct sum of unit inclusions.
\end{definition}

\begin{theorem}[Tensor-product model for the augmented center complex; \ClaimStatusProvedHere]
% label removed: thm:w-augmented-center-tensor-model
Let $\mathfrak N$ be a decoupled boundary simplex, and for each $i\in I$ set
\[
K_i:=\bigl[\Bbbk\xrightarrow{\eta_i} \Bbbk\mathbf{1}\oplus \widetilde Z_i\bigr],
\]
concentrated in cohomological degrees $0$ and $1$, where $\eta_i(1)=\mathbf{1}$.
Then there is a canonical isomorphism of cochain complexes
\[
\widetilde C_Z^\bullet(\mathfrak N)
\xrightarrow{\sim}
\left(\widehat\bigotimes_{i=1}^m K_i\right)[-1].
\]
Consequently,
\[
H^p\bigl(C_Z^\bullet(\mathfrak N)\bigr)=0
\quad (0\le p<m-1),
\qquad
H^{m-1}\bigl(C_Z^\bullet(\mathfrak N)\bigr)
\cong
\widehat\bigotimes_{i=1}^m \widetilde Z_i.
\]
\end{theorem}

\begin{proof}
Expand the completed tensor product $\widehat\bigotimes_i K_i$. A pure tensor lies in cohomological
degree $r$ exactly when $r$ of the factors lie in $\Bbbk\mathbf{1}\oplus \widetilde Z_i$ and the remaining
$m-r$ lie in $\Bbbk$. Such a choice is equivalent to a subset $S\subseteq I$ of cardinality $r$, and the
corresponding summand is canonically
\[
\widehat\bigotimes_{i\in S}(\Bbbk\mathbf{1}\oplus \widetilde Z_i)
\cong Z_S.
\]
Thus the degree-$r$ part of $\widehat\bigotimes_i K_i$ identifies with
\[
\bigoplus_{|S|=r} Z_S = C_Z^{r-1}(\mathfrak N),
\]
which explains the shift by $[-1]$.

Under this identification, the tensor-product differential inserts the unit in one additional factor, with
the usual Koszul sign determined by the position of that factor. This is exactly the alternating sum in the
augmented \v{C}ech differential of Definition~\ref{def:w-augmented-center-complex}. Hence the two complexes
are canonically isomorphic.

Each $K_i$ has $H^0(K_i)=0$ and $H^1(K_i)\cong \widetilde Z_i$. The completed K"unneth theorem therefore gives
\[
H^r\left(\widehat\bigotimes_{i=1}^m K_i\right)=0\quad (r\neq m),
\qquad
H^m\left(\widehat\bigotimes_{i=1}^m K_i\right)
\cong
\widehat\bigotimes_{i=1}^m \widetilde Z_i.
\]
After shifting by $[-1]$, the cohomology of $\widetilde C_Z^\bullet(\mathfrak N)$ is concentrated in degree $m-1$.
Since $\widetilde C_Z^\bullet$ and $C_Z^\bullet$ agree in nonnegative degrees, the same statement holds for
$C_Z^\bullet$.
\end{proof}

\begin{corollary}[Pure $m$-body coupling theorem; \ClaimStatusProvedHere]
% label removed: cor:w-pure-mbody-coupling-theorem
Let $\mathfrak N$ be a decoupled boundary simplex of order $m$. Then
\[
T^{1,[m]}_{X,\mathrm{pure}}(\mathfrak N)
\cong
\left(\widehat\bigotimes_{i=1}^m \widetilde Z_i\right)^\vee\otimes \omega_X.
\]
Thus pure $m$-body first-order couplings are controlled by the full tensor product of the reduced
facewise derived centers.
\end{corollary}

\begin{proof}
Combine Theorem~\ref{thm:w-boundary-cech-duality} with
Theorem~\ref{thm:w-augmented-center-tensor-model}.
\end{proof}

\begin{theorem}[Higher Mayer--Vietoris descent for couplings; \ClaimStatusProvedHere]
% label removed: thm:w-higher-mayer-vietoris-couplings
Let $\mathfrak N$ be any boundary simplex of order $m$.
\begin{enumerate}[label=(\roman*)]
\item The obstruction to descending a compatible family of codimension-$1$ couplings
\[
\{\lambda_S\in T_S\}_{|S|=m-1}
\]
to a coupling on the top face $I$ is the cohomology class
\[
[\lambda]\in H^1\bigl(C_T^\bullet(\mathfrak N)\bigr)
\cong
H^{m-2}\bigl(C_Z^\bullet(\mathfrak N)\bigr)^\vee\otimes \omega_X.
\]
\item If $H^{m-2}(C_Z^\bullet(\mathfrak N))=0$, then every compatible codimension-$1$ family descends to
some coupling on the top face.
\item If, in addition, $H^{m-1}(C_Z^\bullet(\mathfrak N))=0$, then such a descent is unique.
\end{enumerate}
More generally, $H^q(C_T^\bullet(\mathfrak N))$ is the obstruction space for descending compatible
codimension-$q$ face data one step upward.
\end{theorem}

\begin{proof}
A codimension-$1$ family is by definition a degree-$1$ cochain in $C_T^\bullet(\mathfrak N)$. It is
compatible on codimension-$2$ faces exactly when it is a cocycle. Such a cocycle descends from a top-face
coupling precisely when it lies in the image of
\[
\partial\colon C_T^0(\mathfrak N)\to C_T^1(\mathfrak N),
\]
so the obstruction is its class in $H^1(C_T^\bullet)$. This proves (i), and (ii) is the vanishing criterion
$H^1(C_T^\bullet)=0$ rewritten via Theorem~\ref{thm:w-boundary-cech-duality}. If a descent exists, any two
choices differ by a degree-$0$ cocycle, i.e. by an element of $H^0(C_T^\bullet)$; thus the descent is unique
when $H^0(C_T^\bullet)=0$, which is equivalent to $H^{m-1}(C_Z^\bullet)=0$ by Theorem~\ref{thm:w-boundary-cech-duality}.
The final statement is the same argument in degree $q$.
\end{proof}

\begin{definition}[Filtered boundary simplex]
% label removed: def:w-filtered-boundary-simplex
A boundary simplex $\mathfrak N$ is \emph{filtered} if each $Z_S$ carries a complete descending filtration
$F^\bullet Z_S$ such that every face map $(\rho_{S,T})_*$ preserves filtration.
\end{definition}

\begin{theorem}[Filtered stability of pure $m$-body couplings; \ClaimStatusProvedHere]
% label removed: thm:w-filtered-stability-pure-mbody
Let $\mathfrak N$ be a filtered boundary simplex. Assume its associated graded simplex
$\operatorname{gr}\mathfrak N$ is decoupled in the sense of Definition~\ref{def:w-decoupled-simplex}, with reduced
centers $\widetilde Z_1,\dots,\widetilde Z_m$. Then:
\begin{enumerate}[label=(\alph*)]
\item the filtered complex $C_Z^\bullet(\mathfrak N)$ has a convergent spectral sequence whose $E_1$-page is
\[
E_1^p \cong H^p\bigl(C_Z^\bullet(\operatorname{gr}\mathfrak N)\bigr);
\]
\item this spectral sequence collapses at $E_1$;
\item the top cohomology inherits a canonical filtration with
\[
\operatorname{gr}H^{m-1}\bigl(C_Z^\bullet(\mathfrak N)\bigr)
\cong
\widehat\bigotimes_{i=1}^m \widetilde Z_i,
\]
while $H^p(C_Z^\bullet(\mathfrak N))=0$ for $0\le p<m-1$;
\item therefore the pure $m$-body coupling space carries a canonical filtration satisfying
\[
\operatorname{gr}T^{1,[m]}_{X,\mathrm{pure}}(\mathfrak N)
\cong
\left(\widehat\bigotimes_{i=1}^m \widetilde Z_i\right)^\vee\otimes \omega_X.
\]
\end{enumerate}
\end{theorem}

\begin{proof}
The facewise filtrations on the $Z_S$ induce a complete filtration on the direct sums
$C_Z^p(\mathfrak N)$, hence on the whole \v{C}ech complex. Standard filtered-complex theory gives a convergent
spectral sequence with
\[
E_1^p\cong H^p\bigl(C_Z^\bullet(\operatorname{gr}\mathfrak N)\bigr).
\]
By assumption, the associated graded simplex is decoupled, so Theorem~\ref{thm:w-augmented-center-tensor-model}
shows that $E_1^p=0$ for $p<m-1$ and
\[
E_1^{m-1}\cong \widehat\bigotimes_{i=1}^m \widetilde Z_i.
\]
Since the $E_1$-page is concentrated in a single cohomological degree, all higher differentials vanish and the
spectral sequence collapses at $E_1$. Therefore the final cohomology is also concentrated in degree $m-1$, and its
associated graded is the stated tensor product. The final assertion follows from Theorem~\ref{thm:w-boundary-cech-duality}.
\end{proof}

\begin{remark}[Meaning of filtered stability]
% label removed: rem:w-filtered-stability-meaning
The theorem shows that, once the facewise derived centers admit a compatible filtration with decoupled associated
graded, the entire first-order pure $m$-body coupling theory is a \emph{filtered deformation} of the quadratic tensor-product law.
In particular, the non-quadratic $\mathcal W$-sector does not destroy pure higher-body couplings at first order; it refines them.
\end{remark}
% ===== End {fn} =====


% ===== Begin ym_instanton_ir_extension.tex =====
% New frontier package: instanton completion, screening, confinement, and mass-gap bridge.
% Suggested insertion point: after the higher-body \v{C}ech package for non-quadratic W-boundaries,
% or as a new subsection near the chapter's future-directions block.

\subsection{Instanton completion, screening, confinement, and the untwisted infrared bridge}
% label removed: subsec:instanton-screening-ir-bridge

\begin{summary}[New first-principles ingredients]
The preceding boundary packages organize perturbative and twisted Yang--Mills sectors.
To bring instantons, confinement, and the untwisted mass-gap problem into scope, three
new ingredients are required:
\begin{enumerate}[label=(\roman*)]
\item an \emph{instanton-completed} version of the bar/cobar and deformation packages, built from
Novikov-type completions and compatible with Koszul duality;
\item a \emph{central screening formalism} on the dual side, extracting the genuinely visible
infrared center by quotienting by instanton-generated screening charges;
\item a \emph{Hilbert-screening bridge} that translates visible-center collapse into a spectral-gap
criterion for an untwisted Hamiltonian whenever one can dominate the Hamiltonian by the screening
Laplacian.
\end{enumerate}
The results of this subsection do \emph{not} solve the Yang--Mills mass-gap problem. Rather, they
materially enlarge the framework so that instanton sectors, confinement criteria, and a conditional
mass-gap mechanism become internal objects of the same formalism.
\end{summary}

\begin{definition}[Instanton charge monoid and Novikov ring]
% label removed: def:instanton-charge-monoid
An \emph{instanton charge monoid} is a commutative, cancellative, pointed monoid
$(\Gamma,+,0)$ equipped with a proper additive height function
\[
E\colon \Gamma \longrightarrow \mathbb Z_{\ge 0}
\]
(i.e. $E(\beta+\gamma)=E(\beta)+E(\gamma)$ and every fiber $E^{-1}([0,N])$ is finite).
Its \emph{Novikov ring} is
\[
\Lambda_{\Gamma}
:=
\left\{\sum_{\beta\in\Gamma} c_{\beta} q^{\beta}
\;\middle|\;
\forall N,\; \#\{\beta\mid c_{\beta}\neq 0,\;E(\beta)\le N\}<\infty
\right\}.
\]
The $q$-adic filtration is
\[
F^{\ge N}\Lambda_{\Gamma}
:=
\left\{\sum c_{\beta}q^{\beta}\;\middle|\; c_{\beta}=0\text{ whenever }E(\beta)<N\right\}.
\]
\end{definition}

\begin{definition}[Instanton-graded homotopy algebra]
% label removed: def:instanton-graded-homotopy-algebra
Let $M=\bigoplus_{\beta\in\Gamma} M_{\beta}$ be a cohomologically graded vector space.
A complete filtered $A_{\infty}$-structure on $M$ is called \emph{instanton-graded} if:
\begin{enumerate}[label=(\alph*)]
\item each higher multiplication preserves total charge,
\[
m_n\bigl(M_{\beta_1},\dots,M_{\beta_n}\bigr)
\subseteq M_{\beta_1+\cdots+\beta_n};
\]
\item each charge slice $M_{\beta}$ is finite-dimensional;
\item for every $N$, only finitely many charge decompositions
$\beta_1+\cdots+\beta_n=\beta$ occur with $E(\beta)\le N$.
\end{enumerate}
Its \emph{Novikov completion} is
\[
M^{\Lambda}:=\prod_{\beta\in\Gamma} M_{\beta}\,q^{\beta}
\cong M\widehat\otimes \Lambda_{\Gamma}.
\]
\end{definition}

\begin{theorem}[Novikov completion theorem; \ClaimStatusProvedHere]
% label removed: thm:novikov-completion-theorem
Let $M$ be an instanton-graded complete filtered $A_{\infty}$-algebra.
Then:
\begin{enumerate}[label=(\alph*)]
\item the formulas
\[
m_n^{\Lambda}(x_1q^{\beta_1},\dots,x_nq^{\beta_n})
:=m_n(x_1,\dots,x_n)q^{\beta_1+\cdots+\beta_n}
\]
extend uniquely by continuity to a complete filtered $A_{\infty}$-structure on $M^{\Lambda}$;
\item the continuous Hochschild cochain complex decomposes chargewise:
\[
CC^{\bullet}_{\mathrm{cont}}(M^{\Lambda})
\cong
\prod_{\beta\in\Gamma} CC^{\bullet}(M)_{\beta}\, q^{\beta};
\]
\item the derived center completes accordingly:
\[
Z_{\infty}(M^{\Lambda})
:=HH^0\bigl(M^{\Lambda}\bigr)
\cong
\prod_{\beta\in\Gamma} HH^0(M)_{\beta}\, q^{\beta}.
\]
\end{enumerate}
\end{theorem}

\begin{proof}
Part (a): the defining formulas are compatible with the $q$-adic filtration because the total charge is
additive. By Definition~\ref{def:instanton-graded-homotopy-algebra}(c), each coefficient of a fixed total
charge in $m_n^{\Lambda}(x_1,\dots,x_n)$ is a finite sum, so the operations are well-defined and continuous.
The Stasheff identities hold coefficientwise in every charge slice, hence on the completion.

Part (b): a continuous cochain on $M^{\Lambda}$ is determined by its components on fixed total input
and output charge. By finite-dimensionality of each charge slice and the $q$-adic completeness, continuity
is equivalent to giving a compatible family of ordinary Hochschild cochains on the finite-charge truncations.
Taking the inverse limit over the truncations yields the displayed product decomposition.

Part (c) is degree-zero cohomology of part~(b).
\end{proof}

\begin{definition}[Instanton-completed boundary datum]
% label removed: def:instanton-completed-boundary-datum
Let $(X,\mathcal T,\mathcal B)$ be a Koszul-admissible twisted Yang--Mills boundary datum in the sense of
Definition~\ref{def:twisted-ym-boundary-datum}. An \emph{instanton-completed lift} of $(X,\mathcal T,\mathcal B)$
is a choice of:
\begin{enumerate}[label=(\roman*)]
\item an instanton charge monoid $\Gamma$;
\item a $\Gamma$-grading on the boundary chiral algebra $\cA_{\mathcal B}$ and on a transferred dual
$M_{\mathcal B}$ compatible with all operations;
\item finite-type hypotheses sufficient to apply Theorem~\ref{thm:novikov-completion-theorem}
chargewise.
\end{enumerate}
We then write
\[
\widehat{\cA}_{\mathcal B}:=\cA_{\mathcal B}^{\Lambda},
\qquad
\widehat M_{\mathcal B}:=M_{\mathcal B}^{\Lambda}.
\]
\end{definition}

\begin{theorem}[Instanton-completed tangent-to-center principle; \ClaimStatusProvedHere]
% label removed: thm:instanton-completed-tangent-center
Let $(X,\mathcal T,\mathcal B)$ be an instanton-completed boundary datum.
Assume the ordinary tangent-to-center principle holds chargewise for the underlying dual pair.
Then there is a canonical isomorphism
\[
HH^2_{\mathrm{chiral}}(\widehat{\cA}_{\mathcal B})
\xrightarrow{\ \sim\ }
Z_{\infty}(\widehat M_{\mathcal B})^{\vee}\otimes \omega_X.
\]
Equivalently, first-order boundary deformations in the instanton-completed theory are dual to the
instanton-completed derived center.
\end{theorem}

\begin{proof}
By Theorem~\ref{thm:novikov-completion-theorem}, both Hochschild and center data decompose by total
charge and complete $q$-adically. Applying Theorem~\ref{thm:twisted-ym-tangent-center} or
Theorem~\ref{thm:twisted-ym-tangent-factors-through-center} on each charge slice gives
\[
HH^2_{\mathrm{chiral}}\bigl((\cA_{\mathcal B})_{\beta}\bigr)
\cong HH^0\bigl((M_{\mathcal B})_{\beta}\bigr)^{\vee}\otimes \omega_X.
\]
Taking the completed product over $\beta\in\Gamma$ yields the result.
\end{proof}

\begin{definition}[Central screening sequence]
% label removed: def:central-screening-sequence
Let $M$ be an instanton-completed transferred dual.
A \emph{central screening sequence} on $M$ is a finite tuple
\[
\mathbf s=(s_1,\dots,s_r),\qquad s_i\in Z_{\infty}(M),
\]
of degree-zero central classes such that each $s_i$ has strictly positive instanton valuation:
\[
v_q(s_i):=\min\{E(\beta)\mid (s_i)_{\beta}\neq 0\}>0.
\]
The ideal generated by $\mathbf s$ is written
\[
I_{\mathbf s}:=(s_1,\dots,s_r)\subset Z_{\infty}(M).
\]
\end{definition}

\begin{definition}[Visible center and visible first-order couplings]
% label removed: def:visible-center-and-couplings
Let $(X,\mathcal T,\mathcal B;\mathbf s)$ be an instanton-completed boundary datum together with a central
screening sequence on $\widehat M_{\mathcal B}$.
Define the \emph{visible center}
\[
Z_{\mathrm{vis}}(\mathcal B;\mathbf s)
:=
Z_{\infty}(\widehat M_{\mathcal B})/I_{\mathbf s}.
\]
Via Theorem~\ref{thm:instanton-completed-tangent-center}, define the space of
\emph{visible first-order couplings} by
\[
T^1_{X,\mathrm{vis}}(\mathcal B;\mathbf s)
:=
\bigl\{\lambda\in HH^2_{\mathrm{chiral}}(\widehat{\cA}_{\mathcal B})\;\big|\; \lambda(I_{\mathbf s})=0\bigr\}.
\]
\end{definition}

\begin{theorem}[Screened tangent-to-center principle; \ClaimStatusProvedHere]
% label removed: thm:screened-tangent-center
For every instanton-completed boundary datum equipped with a central screening sequence,
\[
T^1_{X,\mathrm{vis}}(\mathcal B;\mathbf s)
\xrightarrow{\ \sim\ }
Z_{\mathrm{vis}}(\mathcal B;\mathbf s)^{\vee}\otimes \omega_X.
\]
In particular, central screenings remove precisely those first-order boundary couplings that pair
nontrivially with the screening ideal.
\end{theorem}

\begin{proof}
Theorem~\ref{thm:instanton-completed-tangent-center} identifies
\[
HH^2_{\mathrm{chiral}}(\widehat{\cA}_{\mathcal B})
\cong Z_{\infty}(\widehat M_{\mathcal B})^{\vee}\otimes \omega_X.
\]
Under this identification, the condition $\lambda(I_{\mathbf s})=0$ says exactly that $\lambda$ factors through the quotient
$Z_{\infty}(\widehat M_{\mathcal B})\twoheadrightarrow Z_{\mathrm{vis}}(\mathcal B;\mathbf s)$. Dualizing the quotient map
produces the stated isomorphism.
\end{proof}

\begin{definition}[Central Koszul screening complex]
% label removed: def:central-koszul-screening-complex
For a commutative algebra $R$ and a sequence $\mathbf s=(s_1,\dots,s_r)\subset R$, define the
\emph{central Koszul screening complex}
\[
\Koszul^{\bullet}(R;\mathbf s)
:=
\bigl(R\otimes \Lambda^{\bullet}(\epsilon_1,\dots,\epsilon_r),\ d_{\mathbf s}=\sum_{i=1}^r \epsilon_i s_i\bigr),
\]
with $|\epsilon_i|=-1$.
For an instanton-completed boundary datum we abbreviate
\[
\Koszul^{\bullet}_{\mathrm{vis}}(\mathcal B;\mathbf s)
:=
\Koszul^{\bullet}\bigl(Z_{\infty}(\widehat M_{\mathcal B});\mathbf s\bigr).
\]
\end{definition}

\begin{proposition}[Koszul identification of the visible center; \ClaimStatusProvedHere]
% label removed: prop:koszul-identification-visible-center
If $\mathbf s$ is a regular sequence on $Z_{\infty}(\widehat M_{\mathcal B})$, then
\[
H^0\bigl(\Koszul^{\bullet}_{\mathrm{vis}}(\mathcal B;\mathbf s)\bigr)
\cong Z_{\mathrm{vis}}(\mathcal B;\mathbf s),
\qquad
H^{-p}\bigl(\Koszul^{\bullet}_{\mathrm{vis}}(\mathcal B;\mathbf s)\bigr)=0\quad (p>0).
\]
\end{proposition}

\begin{proof}
This is the standard exactness of the Koszul complex for a regular sequence, applied to the commutative
algebra $Z_{\infty}(\widehat M_{\mathcal B})$.
\end{proof}

\begin{definition}[Instanton filtration and initial screening system]
% label removed: def:initial-screening-system
Let $R=Z_{\infty}(\widehat M_{\mathcal B})$ with its $q$-adic filtration.
For $x\in R$, write $\operatorname{in}(x)\in \operatorname{gr}R$ for the initial term with respect to the
instanton valuation. For a central screening sequence $\mathbf s=(s_1,\dots,s_r)$, set
\[
\operatorname{in}(\mathbf s):=\bigl(\operatorname{in}(s_1),\dots,\operatorname{in}(s_r)\bigr).
\]
\end{definition}

\begin{theorem}[Instanton-screening spectral sequence; \ClaimStatusProvedHere]
% label removed: thm:instanton-screening-spectral-sequence
The $q$-adic filtration on $R$ induces a convergent spectral sequence
\[
E_1^{-p,*}
\cong
H^{-p}\Bigl(\Koszul^{\bullet}\bigl(\operatorname{gr}R;\operatorname{in}(\mathbf s)\bigr)\Bigr)
\Longrightarrow
H^{-p}\Bigl(\Koszul^{\bullet}(R;\mathbf s)\Bigr).
\]
In particular, if the initial screening sequence $\operatorname{in}(\mathbf s)$ is regular and
\[
\operatorname{gr}R/\bigl(\operatorname{in}(s_1),\dots,\operatorname{in}(s_r)\bigr)\cong \mathbb C,
\]
then
\[
Z_{\mathrm{vis}}(\mathcal B;\mathbf s)\cong \mathbb C.
\]
\end{theorem}

\begin{proof}
The filtered complex is $\Koszul^{\bullet}(R;\mathbf s)$ with the filtration induced from the $q$-adic
filtration on $R$. Since each $d_{\mathbf s}$ raises instanton valuation by at least $v_q(s_i)>0$, the
associated graded differential is precisely the Koszul differential defined by the initial terms
$\operatorname{in}(s_i)$. Completeness and separatedness of the $q$-adic filtration yield convergence.
The final statement follows from Proposition~\ref{prop:koszul-identification-visible-center} applied on the
$E_1$-page together with collapse of the spectral sequence.
\end{proof}

\begin{definition}[Visible defect category]
% label removed: def:visible-defect-category
Let $\mathcal A_{\mathrm{def}}:=H^0(\widehat M_{\mathcal B})$, and let $S\subset Z(\mathcal A_{\mathrm{def}})$ be the
multiplicative set generated by a central screening sequence $\mathbf s$.
Assume $\mathcal A_{\mathrm{def}}$ is left noetherian.
Define the \emph{visible defect category}
\[
\mathsf{Def}_{\mathrm{vis}}(\mathcal B;\mathbf s)
:=S^{-1}\bigl(\mathcal A_{\mathrm{def}}\text{-}\mathrm{mod}_{\mathrm{fg}}\bigr),
\]
namely the exact central localization of finitely generated $\mathcal A_{\mathrm{def}}$-modules.
A finitely generated $\mathcal A_{\mathrm{def}}$-module $N$ is called \emph{screened} if every element of $N$
is annihilated by some element of $S$.
\end{definition}

\begin{theorem}[Central localization criterion for confinement; \ClaimStatusProvedHere]
% label removed: thm:central-localization-confinement
Assume the hypotheses of Definition~\ref{def:visible-defect-category}.
Then:
\begin{enumerate}[label=(\alph*)]
\item a finitely generated module maps to zero in $\mathsf{Def}_{\mathrm{vis}}(\mathcal B;\mathbf s)$ if and only if it is screened;
\item if every nontrivial simple finitely generated $\mathcal A_{\mathrm{def}}$-module is screened and
$Z_{\mathrm{vis}}(\mathcal B;\mathbf s)\cong \mathbb C$, then $\mathsf{Def}_{\mathrm{vis}}(\mathcal B;\mathbf s)$ has a unique simple object,
namely the vacuum module.
\end{enumerate}
\end{theorem}

\begin{proof}
Part (a) is the standard characterization of the kernel of central localization: a module becomes zero after
inverting $S$ exactly when every element is killed by some element of $S$.

For part (b), all nontrivial simple objects are screened by assumption, hence vanish in the localization.
The vacuum module survives because its central character factors through the scalar quotient
$Z_{\mathrm{vis}}\cong \mathbb C$. Since no other simple object survives, the localized category has a unique
simple object.
\end{proof}

\begin{remark}[Physical interpretation of Theorem~\ref{thm:central-localization-confinement}]
% label removed: rem:physical-interpretation-confinement
When line operators are modeled by modules over the dual algebra (as in the perturbative holomorphic-topological
examples discussed in the surrounding literature), Theorem~\ref{thm:central-localization-confinement}
provides a precise mathematical notion of \emph{confinement}: after inverting the instanton-induced screening
charges, only the vacuum asymptotic sector remains visible.
This interpretation is conditional on the physical realization of the defect category, but the localization theorem itself is purely algebraic.
\end{remark}

\begin{definition}[Screened boundary simplex]
% label removed: def:screened-boundary-simplex
Let $\mathfrak N=\{M_S\}_{\emptyset\neq S\subseteq I}$ be a boundary simplex in the sense of
Definition~\ref{def:w-boundary-simplex}.
A \emph{screened boundary simplex} consists of the data of central screening sequences
\[
\mathbf s_S\subset Z_S:=HH^0(M_S)
\qquad (\emptyset\neq S\subseteq I)
\]
such that every face map sends screenings to screenings:
\[
(\rho_{S,T})_*\bigl(\mathbf s_T\bigr)\subseteq \mathbf s_S
\qquad (\emptyset\neq T\subseteq S\subseteq I).
\]
Write
\[
Z_S^{\mathrm{vis}}:=Z_S/(\mathbf s_S),
\qquad
T_S^{\mathrm{vis}}:=(Z_S^{\mathrm{vis}})^{\vee}\otimes \omega_X.
\]
\end{definition}

\begin{theorem}[Higher-body \v{C}ech duality after screening; \ClaimStatusProvedHere]
% label removed: thm:screened-cech-duality
Let $\mathfrak N$ be a screened boundary simplex of order $m$.
Define
\[
C_{Z,\mathrm{vis}}^p(\mathfrak N):=\bigoplus_{|S|=p+1} Z_S^{\mathrm{vis}},
\qquad
C_{T,\mathrm{vis}}^p(\mathfrak N):=\bigoplus_{|S|=m-p} T_S^{\mathrm{vis}},
\]
with the obvious face maps induced from the simplex structure.
Then for every $0\le q\le m-1$ there is a natural isomorphism
\[
H^q\bigl(C_{T,\mathrm{vis}}^{\bullet}(\mathfrak N)\bigr)
\cong
H^{m-1-q}\bigl(C_{Z,\mathrm{vis}}^{\bullet}(\mathfrak N)\bigr)^{\vee}\otimes \omega_X.
\]
In particular, the pure visible $m$-body coupling space is
\[
T_{X,\mathrm{vis}}^{1,[m]}(\mathfrak N)
:=H^0\bigl(C_{T,\mathrm{vis}}^{\bullet}(\mathfrak N)\bigr)
\cong
H^{m-1}\bigl(C_{Z,\mathrm{vis}}^{\bullet}(\mathfrak N)\bigr)^{\vee}\otimes \omega_X.
\]
\end{theorem}

\begin{proof}
This is the same linear-algebra argument as in the unscreened higher-body \v{C}ech theorem, now applied
facewise to the visible tangent-to-center identification of Theorem~\ref{thm:screened-tangent-center}.
The differentials are dual under the pairing because the screening quotients are compatible with the simplex maps.
\end{proof}

\begin{definition}[Hilbert-screening datum]
% label removed: def:hilbert-screening-datum
A \emph{Hilbert-screening datum} consists of:
\begin{enumerate}[label=(\roman*)]
\item a Hilbert space $\mathcal H$ with vacuum projection $P_0$;
\item a family of pairwise strongly commuting normal operators $S_1,\dots,S_r$ on $\mathcal H$ with common invariant core;
\item the fermionic creation and annihilation operators
\[
\epsilon_i,\iota_i\colon \Lambda^{\bullet}(\mathbb C^r)\to \Lambda^{\bullet}(\mathbb C^r)
\]
satisfying $\epsilon_i\iota_j+\iota_j\epsilon_i=\delta_{ij}$.
\end{enumerate}
The associated \emph{screening differential} on
\[
\mathcal K_{\mathrm{scr}}:=\mathcal H\otimes \Lambda^{\bullet}(\mathbb C^r)
\]
is
\[
\partial_{\mathrm{scr}}:=\sum_{i=1}^r \epsilon_i S_i,
\qquad
\partial_{\mathrm{scr}}^*:=\sum_{i=1}^r \iota_i S_i^*.
\]
Its \emph{screening Laplacian} is
\[
L_{\mathrm{scr}}:=\partial_{\mathrm{scr}}\partial_{\mathrm{scr}}^*+\partial_{\mathrm{scr}}^*\partial_{\mathrm{scr}}.
\]
\end{definition}

\begin{theorem}[Screening Hodge theorem; \ClaimStatusProvedHere]
% label removed: thm:screening-hodge-theorem
For a Hilbert-screening datum,
\[
L_{\mathrm{scr}}=\sum_{i=1}^r S_i^*S_i\otimes \mathrm{id}_{\Lambda^{\bullet}(\mathbb C^r)}.
\]
Consequently,
\[
\ker L_{\mathrm{scr}} = \bigcap_{i=1}^r \ker S_i \otimes \Lambda^{\bullet}(\mathbb C^r),
\]
and the degree-zero harmonic space is exactly the joint kernel of the screenings.
\end{theorem}

\begin{proof}
Using the CAR relations and strong commutativity of the $S_i$, we compute
\[
\partial_{\mathrm{scr}}\partial_{\mathrm{scr}}^*+\partial_{\mathrm{scr}}^*\partial_{\mathrm{scr}}
=
\sum_{i,j}(\epsilon_i\iota_j+\iota_j\epsilon_i)S_iS_j^*
=
\sum_i S_iS_i^*.
\]
Since each $S_i$ is normal, $S_iS_i^*=S_i^*S_i$, giving the stated formula. The kernel statement follows
because a sum of nonnegative commuting operators vanishes exactly when each summand vanishes.
\end{proof}

\begin{corollary}[Screening spectral gap criterion; \ClaimStatusProvedHere]
% label removed: cor:screening-spectral-gap-criterion
Assume there is a constant $\mu>0$ such that
\[
\sum_{i=1}^r S_i^*S_i\ \ge\ \mu^2(1-P_0)
\]
as quadratic forms on $\mathcal H$.
Then
\[
\operatorname{Spec}(L_{\mathrm{scr}})\cap (0,\mu^2)=\varnothing.
\]
Equivalently, the screening Laplacian has a gap of size at least $\mu^2$ above the vacuum.
\end{corollary}

\begin{proof}
By Theorem~\ref{thm:screening-hodge-theorem}, the nonzero spectrum of $L_{\mathrm{scr}}$ is bounded below by the
nonzero spectrum of $\sum_i S_i^*S_i$, hence by $\mu^2$ on the orthogonal complement of the vacuum.
\end{proof}

\begin{definition}[Untwisting bridge datum]
% label removed: def:untwisting-bridge-datum
An \emph{untwisting bridge datum} consists of a Hilbert-screening datum $(\mathcal H,P_0,\{S_i\})$
together with a nonnegative self-adjoint Hamiltonian $H_{\mathrm{YM}}$ on $\mathcal H$ such that the vacuum
is a zero-energy state and the gauge-invariant physical Hilbert space is realized as a closed subspace of $\mathcal H$.
\end{definition}

\begin{theorem}[Conditional mass-gap transfer via screening domination; \ClaimStatusConditional]
% label removed: thm:conditional-mass-gap-transfer
Let $(\mathcal H,P_0,\{S_i\},H_{\mathrm{YM}})$ be an untwisting bridge datum.
Assume:
\begin{enumerate}[label=(\alph*)]
\item the screening positivity hypothesis of Corollary~\ref{cor:screening-spectral-gap-criterion} holds with constant $\mu>0$;
\item on the orthogonal complement of the vacuum,
\[
H_{\mathrm{YM}}\ \ge\ \alpha\, L_{\mathrm{scr}}
\]
as quadratic forms for some $\alpha>0$.
\end{enumerate}
Then
\[
\operatorname{Spec}(H_{\mathrm{YM}})\cap (0,\alpha\mu^2)=\varnothing.
\]
In particular, $H_{\mathrm{YM}}$ has a positive mass gap at least $\alpha\mu^2$.
\end{theorem}

\begin{proof}
Restrict to the orthogonal complement of the vacuum. By assumption,
\[
H_{\mathrm{YM}}\ge \alpha L_{\mathrm{scr}} \ge \alpha\mu^2.
\]
The min--max principle therefore excludes spectrum in $(0,\alpha\mu^2)$.
\end{proof}

\begin{corollary}[Stable untwisting under bounded error; \ClaimStatusProvedHere]
% label removed: cor:stable-untwisting-bounded-error
Assume the setup of Theorem~\ref{thm:conditional-mass-gap-transfer}, except that
\[
H_{\mathrm{YM}} = \alpha L_{\mathrm{scr}} + K
\]
with $K$ bounded self-adjoint, $KP_0=0$, and
\[
\|K\|<\alpha\mu^2.
\]
Then
\[
\operatorname{Spec}(H_{\mathrm{YM}})\cap (0,\alpha\mu^2-\|K\|)=\varnothing.
\]
The same conclusion holds uniformly for a norm-continuous family $H_t=\alpha L_{\mathrm{scr}}+K_t$ with
$\sup_t\|K_t\|<\alpha\mu^2$.
\end{corollary}

\begin{proof}
On the orthogonal complement of the vacuum,
\[
H_{\mathrm{YM}}\ge \alpha L_{\mathrm{scr}}-\|K\| \ge (\alpha\mu^2-\|K\|)\,\mathrm{id}.
\]
The family statement is identical, with the same uniform lower bound.
\end{proof}

\subsection{Kazhdan--Lusztig equivalence for the K3 chiral
Hall--Drinfeld double}
\label{subsec:kl-k3-chiral-hd-double}

The boundary-to-Koszul-dual machinery assembled above admits, in the
K3 sub-landscape, a specific Kazhdan--Lusztig realisation of the
chiral quantum group as a small quantum group at a primitive $8$-th
root of unity. The relevant ambient is the $\mathcal{B}$-family of
Vol~III Chapter~\ref{ch:k3-chiral-bialgebra-platonic}, hosted on the
bi-based Ran datum $(\mathrm{Ran}(E^{\mathrm{nod,sm}}_{24}),
\overline{\mathcal{A}_2})$. The Humbert-wall monodromy order is $8$
(Vol~I Chapter~\ref{chap:e3-identification-chain-level-platonic},
Theorem~\ref{thm:humbert-lusztig-mukai-8}); under the Russian-school
Kazhdan--Lusztig equivalence (Duke 1993; J.~AMS~6-7, 1993-1994) applied
to the super-Etingof--Kazhdan Manin-pair quantisation, this fixes a
specific root of unity for the chiral quantum group.

\begin{theorem}[Kazhdan--Lusztig equivalence at $\zeta_8$ for the
K3-BKM chiral bialgebra; \ClaimStatusProvedElsewhere]
\label{thm:kl-k3-bkm-zeta8}
\index{Kazhdan--Lusztig equivalence!K3 BKM}
\index{Humbert divisor!$H_1$}
\index{chiral quantum group!K3 Hall--Drinfeld double}
Let $\mathbf{H}_{\Delta_5}$ be the K3 chiral bialgebra of Vol~III
Chapter~\ref{ch:k3-chiral-bialgebra-platonic}
($=$ super-EK quantisation of the BKM Manin pair
$(\mathfrak{g}_{\Delta_5}, \mathfrak{n}_+^{\mathrm{imag}}\oplus
\mathfrak{h}^{\mathrm{imag,rk\,23}})$ with associator cocycle
$\Phi_{10}/\eta^{24}$). On the Koszul locus of the
$\mathcal{B}$-family, there is an equivalence of
$(\infty,1)$-categories
\[
  \mathrm{Mod}^{\mathrm{ch}}(\mathbf{H}_{\Delta_5})^{\mathrm{loc,KL}}
  \;\simeq\;
  \mathrm{Mod}\bigl(\mathfrak{u}_{q_{\mathrm{K3}}}
    (\widehat{\mathfrak{m}}_{\mathrm{K3}})\bigr),
  \qquad
  q_{\mathrm{K3}} \;=\; e^{2\pi i/8},
\]
where $\mathfrak{u}_{q_{\mathrm{K3}}}$ is Lusztig's small quantum
group on the BKM superalgebra. At this specialisation, $\hbar^2 =
-1/\ell = -1/8$ and the universal identity $\hbar^2 \cdot
K^{\kappa_{\mathrm{ch}}} = -1$ holds with Mukai-doubled Koszul
conductor $K^{\kappa_{\mathrm{ch}}} = 2 c_+(\Lambda_{\mathrm{Muk}}) =
8$.
\end{theorem}

\begin{proof}[Proof (reference)]
Theorem is Corollary~\ref{cor:kl-equivalence-k3-bkm} of Vol~I
Chapter~\ref{chap:e3-identification-chain-level-platonic}. Three
independent verifications (Bruinier 2002 Prop.~5.1 Heegner--Chern
reciprocity, Katzarkov--Pantev--Shepherd-Barron 1998 K3-period local
monodromy, Lusztig 1990 small-quantum-group construction) coincide in
$\mathrm{CH}^1(H_1)\otimes \mathbb{Z}/8$; see Vol~III compute module
\texttt{k3\_yangian\_wave14\_humbert\_monodromy\_8.py}.
\end{proof}

\begin{remark}[Hecke-algebra action and $\Delta_{10}$ Arthur packet]
\label{rem:hecke-action-arthur-ym}
\index{Hecke algebra!$\mathrm{Sp}_4$}
\index{Arthur parameter!$\Delta_{10}$}
The unramified Hecke algebra $\mathcal{H}(\mathrm{Sp}_4(\mathbb{Q}_p),
\mathrm{Sp}_4(\mathbb{Z}_p))$ acts on $S_{10}(\mathrm{Sp}_4
(\mathbb{Z}))$ through its Satake transform, landing in the polynomial
ring on ${}^L\mathrm{Sp}_4 = \mathrm{SO}_5(\mathbb{C})$. The Arthur
parameter
\[
  \psi_{\Delta_{10}}
  \;=\;
  \phi_{\Delta E_6}\,\boxtimes\,\mathrm{Sym}^1
  \colon
  L_F \times \mathrm{SL}_2(\mathbb{C})
  \longrightarrow
  \mathrm{SO}_5(\mathbb{C})
\]
(Ibukiyama--Skoruppa--Zagier, following Andrianov 1987) produces
Hecke eigenvalues at $p$ whose Satake parameters factor through those
of $\Delta E_6\in S_{18}(\mathrm{SL}_2(\mathbb{Z}))$ as
$L_p(s, \Delta_{10}, \mathrm{Spin}) = L_p(s, \Delta E_6)\cdot
\zeta_p(s-9)\cdot \zeta_p(s-8)$. Under
Theorem~\ref{thm:kl-k3-bkm-zeta8}, the Hecke eigenvalue $a_p(\Delta
E_6)/p^{17/2}$ matches the local Casimir spectrum of
$\mathfrak{u}_{q_{\mathrm{K3}}}$ at the prime-$p$ stratum of the small
quantum group; the spin refinement to $\Delta_5$ multiplies by
$\epsilon_p = v_{\Delta_5}(\mathrm{Frob}_p) \in \{\pm 1\}$ (Maass
1964 multiplier system). Cross-references: Vol~I
Theorem~\ref{thm:pref-k3-arthur-parametrisation} (preface, Section 4.7);
Vol~III Chapter~\ref{ch:k3-chiral-bialgebra-platonic}.
\end{remark}

\begin{remark}[Langlands self-duality of the K3-BKM chiral bialgebra]
\label{rem:langlands-self-duality-k3-bkm}
\index{Langlands duality!K3 BKM self-duality}
The K3-BKM is Langlands self-dual in two compatible senses.
\emph{(i) Lattice-arithmetic.} The Cartan lattice $\Lambda^{2,1}_{II}$
is hyperbolic and the BKM Gram matrix symmetric, so the formal
Langlands-dual Cartan is the same lattice.
\emph{(ii) Automorphic.} The L-group of
$\Gamma^{\mathrm{spin}}_{\Delta_5} =
\widehat{\mathrm{Sp}_4(\mathbb{Z})}^{v_{\Delta_5}}$ is the
$\mathrm{Spin}_5 \cong \mathrm{Sp}_4$ double cover of
${}^L\mathrm{Sp}_4 = \mathrm{SO}_5 = \mathrm{PGSp}_4$, recovering
$\mathrm{Sp}_4$ itself (Lorgat 2020, Lemma~1). The Fricke involution
$S\colon (\tau_1, z, \tau_2) \mapsto -(\tau_2, -z, \tau_1)$ on
$\overline{\mathcal{A}_2}$ restricts to an involution on
$\mathbf{H}_{\Delta_5}$ exchanging the two Heisenberg halves of the
BKM and fixing the imaginary rank-$23$ Cartan; $\Phi_{10}$ is
$S$-invariant (Gritsenko--Nikulin 1998). The self-duality is a genuine
mathematical fact about this chiral bialgebra, not a tautology.
\end{remark}

\begin{remark}[What is proved, and what remains]
% label removed: rem:what-is-proved-what-remains-inst
Theorems~\ref{thm:novikov-completion-theorem}--\ref{cor:stable-untwisting-bounded-error} materially enlarge the
mathematical framework.
They do \emph{not} solve untwisted Yang--Mills. The new algebraic core now includes:
\begin{enumerate}[label=(\roman*)]
\item instanton-completed Hochschild and center theory;
\item a precise quotient formalism for screening and confinement;
\item higher-body visible couplings on boundary simplices after screening;
\item a conditional route from screening positivity to an honest mass gap.
\end{enumerate}
What remains genuinely analytic and physical is the construction of a natural untwisting bridge datum
for the actual untwisted four-dimensional Yang--Mills Hamiltonian, together with verification of the
screening domination estimate in Theorem~\ref{thm:conditional-mass-gap-transfer}.
\end{remark}
% ===== End {fn} =====

\section{Synthesis}
\label{sec:yms-wave13-DNA}

\begin{remark}[Yang-Mills class-$\mathcal S$ synthesis on $\Sigma_{0,24}$]
\label{rem:yms-DNA-1}
The the K3 synthesis Yang-Mills class-$\mathcal S$ synthesis at the K3 base point
is class-$\mathcal S$ $A_1$ on $\Sigma_{0,24}$: $4d$ $\mathcal N=2$ theory
with gauge group $A_1^{21}$ (Coulomb rank $21$) coupled at $24$ punctures
with flavour $\mathfrak{su}(2)^{24}$. The $4d$/$2d$ Schur-index chiral
algebra is $\mathbf H_{\Delta_5}$-module category with
$c_{4d} = 107/6, c_{2d} = -214$ (Shapere--Tachikawa $2008$ \S 3
applied to Chacaltana--Distler $2010$ eq.~(5.14) count
$(n_v, n_h) = (63, 88)$; the intermediate $(26, -312)$ retracted
per Vol~I Chapter~\ref{ch:thqg-open-closed-realization} line~1932)
\cite{Gaiotto2009,BeemRastelli2015,ChacaltanaDistler2010}.
Cross-ref Vol.~III Chapter~\ref{ch:k3-chiral-bialgebra-platonic}.
\end{remark}
